\chapter{Review of Literature and Plagiarism Issues}

\section{Historical Background of Malaria Control}

Malaria has plagued humanity for millennia. Descriptions of periodic fevers characteristic of malaria appear in ancient Greek, Roman, Chinese, and Indian medical texts dating back to 2700 BCE. The name "malaria" derives from the Italian "mala aria" meaning "bad air," reflecting the pre-scientific belief that the disease was caused by swamp vapors.

The discovery that mosquitoes transmit malaria by \citet{ross1898} and \citet{grassi1900} revolutionized public health approaches to disease control. Ross demonstrated that malaria parasites develop in mosquitoes, earning the Nobel Prize in Medicine in 1902. This discovery opened the possibility of vector control.

Early control efforts focused on environmental management: draining swamps, filling standing water, and using oil to coat water surfaces to kill mosquito larvae. These approaches were labor-intensive but effective where implemented.

The advent of dichlorodiphenyltrichloroethane (DDT) during World War II marked a turning point in malaria control. DDT was cheap, persistent, and highly effective against adult mosquitoes. Indoor residual spraying (IRS) with DDT became the cornerstone of malaria control.

The Global Malaria Eradication Program (1955-1969), led by the World Health Organization, relied heavily on IRS with DDT. The program achieved remarkable success: malaria was eliminated from Europe, North America, the Caribbean, and parts of Asia. However, the program failed in sub-Saharan Africa due to several factors:

\begin{enumerate}
	\item \textbf{Logistical challenges}: The infrastructure for sustained IRS campaigns was lacking in many African countries.
	
	\item \textbf{Insecticide resistance}: Anopheles mosquitoes developed resistance to DDT and other insecticides.
	
	\item \textbf{Weak health infrastructure}: Many countries lacked the trained personnel and monitoring systems for effective vector control.
	
	\item \textbf{High transmission intensity}: In tropical Africa, transmission intensity was orders of magnitude higher than in areas where elimination succeeded.
\end{enumerate}

After the program ended, malaria resurged dramatically, with incidence and mortality returning to pre-campaign levels or higher. This experience taught the global health community an important lesson: malaria control requires sustained investment; temporary campaigns are not sufficient.

The resurgence of malaria in the 1970s and 1980s prompted a shift toward integrated vector management. Insecticide-treated nets (ITNs) emerged as a promising intervention in the 1990s. The concept was simple: treat bed nets with insecticide to both physically block mosquitoes and kill them on contact.

\section{Evidence for ITN Effectiveness}

\subsection{Randomized Controlled Trials}

The gold standard evidence for ITN effectiveness comes from randomized controlled trials (RCTs). In an RCT, participants are randomly assigned to treatment (ITN) or control (no net or untreated net) groups. Randomization ensures that, in expectation, the groups are balanced on all characteristics, both observed and unobserved.

The first large-scale ITN trial was conducted in The Gambia in the early 1990s. The trial found that ITNs reduced all-cause child mortality by approximately 25\%. This result was so striking that the trial was stopped early by the data safety monitoring board.

Subsequent trials across Africa confirmed these findings. A Cochrane review by \citet{lengeler2004insecticide} synthesized evidence from 22 trials and found that ITNs reduce:

\begin{enumerate}
	\item \textbf{Child mortality}: 17\% reduction (RR 0.83, 95\% CI 0.77-0.89)
	\item \textbf{Episodes of malaria}: 50\% reduction
	\item \textbf{Parasitemia}: 20-30\% reduction
\end{enumerate}

An updated Cochrane review by \citet{pryce2018insecticide} confirmed these findings, reporting a 17\% reduction in \textit{P. falciparum} prevalence (RR 0.83, 95\% CI 0.71-0.98).

The mechanisms underlying these effects are twofold:

\begin{enumerate}
	\item \textbf{Personal protection}: The net physically blocks mosquitoes from biting the person underneath it. Even untreated nets provide some protection, but insecticide treatment substantially enhances efficacy.
	
	\item \textbf{Community protection}: When many people in an area use ITNs, the mosquito population is suppressed. Mosquitoes that attempt to feed are killed by the insecticide before they can transmit parasites. This benefits even non-users in the same area.
\end{enumerate}

\subsection{Observational Studies}

While RCTs provide strong internal validity, their generalizability to real-world settings is limited. RCTs typically:

\begin{enumerate}
	\item Have strict inclusion criteria, excluding many populations
	\item Provide nets free of charge, which may not reflect programmatic distribution
	\item Ensure high coverage through active distribution and follow-up
	\item Re-treat nets regularly, whereas real-world retreatment rates are low
	\item Follow participants for relatively short periods
\end{enumerate}

Observational studies using large-scale survey data have complemented RCT evidence. \citet{lim2011net} conducted a multicountry analysis of Demographic and Health Survey data from 22 African countries, using matching methods to estimate ITN effectiveness. They found a pooled relative reduction in parasitemia prevalence of 20\% (95\% CI 3-35\%) associated with household ITN ownership.

\citet{bhatt2015effect} estimated that ITNs contributed to a 20-30\% reduction in malaria prevalence in Africa between 2000 and 2015. They attributed approximately 68\% of the total reduction in malaria prevalence to ITNs, with the remainder due to artemisinin combination therapies and other interventions.

\subsection{Evidence from Kenya}

Several studies have examined malaria epidemiology in Kenya. \citet{noor2009risks} mapped malaria risk across Kenya using Bayesian geostatistical models, identifying:

\begin{enumerate}
	\item \textbf{High-transmission zones}: Western region near Lake Victoria (Kisumu, Kisii, Kakamega, Bungoma). EIR ranges from 60-300 infective bites per person per year. Malaria is perennial, with peaks during rainy seasons.
	
	\item \textbf{Moderate-transmission zones}: Coastal region (Mombasa, Kilifi, Kwale). Transmission is seasonal, following the bimodal rainfall pattern.
	
	\item \textbf{Low-transmission zones}: Central highlands (Nairobi, Kiambu, Nyeri) and northern arid areas (Turkana, Marsabit). EIR below 10 infective bites per person per year.
\end{enumerate}

\citet{ayele2024assessment} used multilevel logistic regression to analyze malaria transmission in Kenya, demonstrating the importance of accounting for clustering at the enumeration area level. Their ICC estimate of approximately 40\% aligns closely with our finding of 50-58\%.

\subsection{Community-Wide Effects of ITNs}

A critical insight from \citet{hawley2003community} is that ITNs protect not only users but also neighbours by suppressing mosquito populations. This study was conducted in Asembo, western Kenya, as part of a large cluster-randomized trial.

Key findings from Hawley et al. (2003):

\begin{enumerate}
	\item \textbf{Spatial protection}: Children in control compounds within 300 meters of ITN compounds received about the same protection as children actually using nets. For mortality, the hazard ratio was 0.72 (95\% CI 0.60-0.86) for control compounds within 300 meters of ITN compounds.
	
	\item \textbf{Coverage threshold}: This community effect only appeared when coverage exceeded 50\% of nearby compounds. Below 25\% coverage, no effect was observed. This suggests a threshold effect: until coverage reaches a critical level, individual nets provide only personal protection; once coverage exceeds the threshold, community-wide mosquito suppression occurs.
	
	\item \textbf{Mechanism}: The primary mechanism is mosquito mortality. Pyrethroid insecticides kill adult mosquitoes. When enough mosquitoes are killed, the transmission cycle is interrupted. Mosquitoes that attempt to feed in ITN-using compounds die before they can bite again. Over time, the mosquito population declines, and the proportion of infected mosquitoes (sporozoite rate) decreases.
	
	\item \textbf{Policy implications}: High coverage is essential. Low-coverage programs may waste resources without achieving population impact. Mass distribution campaigns must aim for universal coverage (>80\%) to generate community effects.
\end{enumerate}

This finding has profound implications for the interpretation of observational studies. If community effects are present, then comparing ITN users to non-users in the same area underestimates the true effect, because non-users also benefit from the intervention. Conversely, comparing ITN users to non-users in areas without ITNs may overestimate the true effect because non-users in ITN areas benefit from community protection. The DML estimate of 2.8\% reduction likely mixes both personal and community protection, and with cross-sectional data, we cannot separate these mechanisms.

\section{Causal Inference Methods in Epidemiology}

\subsection{Traditional Methods}

Traditional causal inference methods include logistic regression, which estimates associations between exposures and outcomes while adjusting for confounders \citep{austin2011introduction}. The logistic regression model is:

\[
\log\left(\frac{P(Y=1 \mid T, \mathbf{X})}{1-P(Y=1 \mid T, \mathbf{X})}\right) = \beta_0 + \beta_1 T + \boldsymbol{\beta}_2^\top \mathbf{X}
\]

The coefficient $\beta_1$ is interpreted as the log odds ratio for treatment, holding the covariates $\mathbf{X}$ constant. This is a conditional estimate.

Limitations of logistic regression for causal inference:

\begin{enumerate}
	\item \textbf{Linearity assumption}: Assumes the log-odds is linear in the covariates. Non-linear relationships are missed.
	
	\item \textbf{No unmeasured confounding}: Assumes all confounders are measured and included. This is almost certainly false.
	
	\item \textbf{Independence assumption}: Assumes observations are independent. Clustering violates this.
\end{enumerate}

Generalized Linear Mixed Models (GLMM) extend traditional regression by incorporating random effects to account for clustered data structures:

\[
\log\left(\frac{P(Y_{ij}=1 \mid T_{ij}, \mathbf{X}_{ij}, u_j)}{1-P(Y_{ij}=1 \mid T_{ij}, \mathbf{X}_{ij}, u_j)}\right) = \beta_0 + \beta_1 T_{ij} + \boldsymbol{\beta}_2^\top \mathbf{X}_{ij} + u_j
\]

where $u_j \sim N(0, \sigma_u^2)$ is the cluster-specific random intercept. This model estimates cluster-specific (conditional) effects. As \citet{neuhaus1991comparison} demonstrate, when ICC is high, conditional effects differ substantially from marginal effects.

\subsection{Propensity Score Methods}

The propensity score was introduced by \citet{rosenbaum1983central}. It is defined as:

\[
e(\mathbf{X}) = P(T = 1 \mid \mathbf{X})
\]

The key insight is the balancing property:

\[
T \perp \mathbf{X} \mid e(\mathbf{X})
\]

Conditioning on the propensity score balances observed covariates between treated and untreated groups, mimicking a randomized experiment.

\citet{austin2011introduction} provides a comprehensive review of propensity score methods, including:

\begin{enumerate}
	\item \textbf{Matching}: Each treated unit is matched to a control unit with similar propensity score. Unmatched units are discarded. The ATE is the mean difference in outcomes between matched pairs.
	
	\item \textbf{Inverse probability weighting}: Each unit is weighted by the inverse of its probability of receiving the treatment it actually received. This creates a pseudo-population where treatment is independent of covariates.
	
	\item \textbf{Stratification}: Units are divided into strata based on propensity score quintiles. Within each stratum, treatment is approximately randomized. The ATE is the weighted average across strata.
\end{enumerate}

\citet{austin2010optimal} established optimal caliper widths for propensity score matching. The recommended caliper is 0.2 standard deviations of the logit propensity score.

Doubly robust estimation \citep{schafer2008average} combines propensity score weighting with outcome regression:

\[
\hat{\tau}_{\text{DR}} = \frac{1}{n}\sum_{i=1}^{n}\left[\frac{T_i(Y_i - \hat{\mu}_1(\mathbf{X}_i))}{\hat{e}(\mathbf{X}_i)} + \hat{\mu}_1(\mathbf{X}_i)\right] - \frac{1}{n}\sum_{i=1}^{n}\left[\frac{(1-T_i)(Y_i - \hat{\mu}_0(\mathbf{X}_i))}{1 - \hat{e}(\mathbf{X}_i)} + \hat{\mu}_0(\mathbf{X}_i)\right]
\]

The estimator is consistent if either the propensity score model or the outcome model is correctly specified. This double robustness is a major advantage.

\subsection{Machine Learning for Causal Inference}

Double Machine Learning \citep{chernozhukov2018double} combines machine learning algorithms with causal inference techniques. The method uses:

\begin{enumerate}
	\item \textbf{Neyman-orthogonal scores}: The score function is constructed so that its derivative with respect to the nuisance functions is zero at the true parameter. This prevents bias from regularization.
	
	\item \textbf{Cross-fitting}: Nuisance functions are estimated using out-of-sample predictions, preventing overfitting bias.
	
	\item \textbf{Generic machine learning}: Any machine learning algorithm can be used to estimate nuisance functions, as long as it converges at an appropriate rate.
\end{enumerate}

The R-learner \citep{nie2021quasi} estimates heterogeneous treatment effects using orthogonalization. It separates the estimation of the conditional mean outcome and the propensity score from the estimation of treatment effects.

\subsection{The Positivity Assumption}

\citet{westreich2010invited} emphasize the importance of the positivity assumption in causal inference. Positivity requires:

\[
0 < P(T = 1 \mid \mathbf{X}) < 1 \quad \text{for all } \mathbf{X}
\]

This means that for every combination of covariates, there must be both treated and untreated units. If some covariate combination predicts treatment perfectly, the treatment effect is not defined for that combination.

Positivity is often violated when there are deterministic treatment decisions. For example, if the poorest households never have access to ITNs, then positivity fails. In practice, we assess positivity by examining the distribution of propensity scores. If the propensity scores for treated and control units have substantial overlap, positivity is plausible.

\subsection{E-value for Sensitivity Analysis}

\citet{vanderweele2017sensitivity} introduced the E-value, which quantifies the minimum strength of association that an unmeasured confounder would need to fully explain away the observed effect. \citet{chung2023evalue} provide an accessible introduction.

For a protective effect with risk ratio $RR < 1$, the E-value is:

\[
\text{E-value} = \frac{1}{RR} + \sqrt{\frac{1}{RR} \times \left(\frac{1}{RR} - 1\right)}
\]

Interpretation:

\begin{enumerate}
	\item $\text{E-value} = 1$: No unmeasured confounding needed to explain away the effect.
	\item $\text{E-value} = 2$: Unmeasured confounder would need risk ratio of at least 2 with both treatment and outcome.
	\item $\text{E-value} = 3$: Unmeasured confounder would need risk ratio of at least 3 with both treatment and outcome.
\end{enumerate}

Larger E-values indicate greater robustness. For example, an E-value of 3.05 (our result) means that an unmeasured confounder would need to be three times more likely among ITN users and three times more likely among malaria cases to completely explain away the observed effect. This is quite strong, but not impossible.

\section{Gaps in the Literature}

Despite extensive research, several gaps remain:

\begin{enumerate}
	\item \textbf{Lack of causal framing}: Most observational studies report associations rather than effects estimated under explicit causal assumptions. This study explicitly states assumptions and, where possible, tests them.
	
	\item \textbf{Ignoring clustering}: Many analyses ignore the hierarchical structure of survey data, despite evidence of substantial clustering (ICC = 50-58\% in Kenya). This study uses GLMM and cluster-aware cross-fitting.
	
	\item \textbf{Single estimand}: Few studies distinguish between conditional and marginal effects, which have different interpretations for policy. This study estimates both.
	
	\item \textbf{Parametric assumptions}: Traditional methods assume linearity. This study uses flexible machine learning learners.
	
	\item \textbf{No sensitivity analysis}: Few studies assess robustness to unmeasured confounding. This study computes E-values.
\end{enumerate}
